\documentclass{article}
\usepackage[utf8]{inputenc}
\usepackage{a4wide}

\begin{document}

\section*{Startup}

หลังจากจบค่าย สอวน ปฐมาภรณ์ได้คิดการใหญ่ อยากเปิดบริษัท startup

ที่ทำงานเกี่ยวกับการสอนติว สอวน. ออนไลน์ ปฐมาภรณ์ครุ่นคิดเกี่ยวกับการเปิดบริษัทมาก

ทั้งเรื่องงบประมาณ ต้นทุนที่ต้องใช้ กำไร รวมถึงพนักงานด้วย

เนื่องจากปฐมาภรณ์เองก็ยังเด็กจึงมีแนวคิดในการจัดการโครงสร้างของบริษัทให้ตนเองบริหารงานได้ง่ายเพื่อที่จะไม่รบกวนเวลาเรียนของตนเองมากนัก

จากแผนที่วางไว้ว่าจะมีพนักงานทั้งหมด n คน โครงสร้างการบริหารจะใช้โครงสร้างแบบลำดับชั้น

"หัวหน้างาน-ลูกน้อง" (หมายความว่า\ul{\textbf{พนักงานแต่ละคนจะมีหัวหน้างานแค่คนเดียว

ยกเว้นคนหนึ่งที่ไม่มีหัวหน้างาน}})



ในการจัดหัวหน้างาน และลูกน้อง

ปฐมาภรณ์ใช้วิธีให้พนักงานที่อยากจะเป็นหัวหน้างานให้กับพนักงานคนอื่นๆ ยื่นใบสมัคร

รายละเอียดของใบสมัครแต่ละใบจะระบุว่า "พนักงาน a{[}i{]}

พร้อมที่จะเป็นหัวหน้างานให้พนักงาน b{[}i{]} โดยมีค่าใช้จ่ายเพิ่มเติม c{[}i{]}

หากได้รับการแต่งตั้งตามเงื่อนไขนี้" ซึ่งการสมัครนี้จะมีเงื่อนไขว่า ระดับคุณสมบัติ

(qualification) หรือความสามารถของพนักงาน a{[}i{]}

จะต้องมากกว่าระดับคุณสมบัติหรือความสามารถของพนักงาน b{[}i{]} เสมอ

ซึ่งในรอบการสมัครมีใบสมัครส่งมาทั้งสิ้น m ใบ



\ul{งานของคุณ}: ช่วยปฐมาภรณ์คำนวณหาค่าใช้จ่ายที่ต่ำที่สุดในการสร้างโครงสร้างลำดับชั้น

``หัวหน้างาน-ลูกน้อง'' นี้ หรือหากไม่สามารถสร้างโครงสร้างนี้ได้ ก็ให้บอกว่าเป็นไปไม่ได้



\hfill\break



{\def\LTcaptype{none} % do not increment counter

\begin{longtable}[]{@{}

  >{\raggedright\arraybackslash}p{(\linewidth - 4\tabcolsep) * \real{0.3333}}

  >{\raggedright\arraybackslash}p{(\linewidth - 4\tabcolsep) * \real{0.3333}}

  >{\raggedright\arraybackslash}p{(\linewidth - 4\tabcolsep) * \real{0.3333}}@{}}

\toprule\noalign{}

\begin{minipage}[b]{\linewidth}\centering

ข้อมูลเข้า

\end{minipage} & \begin{minipage}[b]{\linewidth}\centering

ข้อมูลออก

\end{minipage} & \begin{minipage}[b]{\linewidth}\centering

คำอธิบาย

\end{minipage} \\

\midrule\noalign{}

\endhead

\bottomrule\noalign{}

\endlastfoot

\begin{minipage}[t]{\linewidth}\raggedright

3\\

9 4 5\\

5\\

3 2 4\\

1 2 4\\

3 2 8\\

1 3 5\\

3 2 5\\

\strut

\end{minipage} & 9 & \begin{minipage}[t]{\linewidth}\raggedright

ในตัวอย่างนี้หนึ่งในวิธีที่เป็นไปได้ในการสร้างลำดับชั้นคือนำใบสมัครที่มีหมายเลข 1 และ 4 มาใช้

ซึ่งจะทำให้ค่าใช้จ่ายรวมเป็น 9\\

\strut

\end{minipage} \\

\begin{minipage}[t]{\linewidth}\raggedright

3\\

1 3 5\\

2\\

3 1 3\\

3 1 1\\

\strut

\end{minipage} & \begin{minipage}[t]{\linewidth}\raggedright

-1\\

\strut

\end{minipage} & \begin{minipage}[t]{\linewidth}\raggedright

ในกรณีนี้ไม่สามารถสร้างลำดับชั้นที่ถูกต้องได้\\

เนื่องจากพนักงานที่มีคุณสมบัติสูงสุด (พนักงาน 3) ต้องเป็นหัวหน้าของพนักงาน 1\\

แต่ใบสมัครที่มีอยู่ไม่สามารถทำให้เกิดลำดับชั้นที่ถูกต้องได้ ดังนั้นคำตอบคือ -1\\

\strut

\end{minipage} \\

\end{longtable}

}



\subsection*{Input}
บรรทัดแรก จำนวนเต็ม n (1 ≤ n ≤ 1000) --- หมายถึงจำนวนพนักงานในบริษัท



บรรทัดถัดไปประกอบด้วยตัวเลข n ตัว คั่นด้วยช่องว่าง --- หมายถึงระดับคุณสมบัติ

(qualification) ของแต่ละพนักงาน (0 ≤ q{[}j{]} ≤ 10\^{}6)



บรรทัดถัดไป จำนวน m (0 ≤ m ≤ 10000) --- หมายถึงจำนวนใบสมัครที่ได้รับ



จากนั้นมี m บรรทัดที่แต่ละบรรทัดประกอบด้วยตัวเลขสามตัว a{[}i{]}, b{[}i{]} และ

c{[}i{]} (1~ ≤~ a{[}i{]}, b{[}i{]}~ ≤~ n, 0 ≤ c{[}i{]} ≤ 10\^{}6) ---

หมายความว่า พนักงาน a{[}i{]} ยินดีที่จะเป็นหัวหน้างานให้พนักงาน b{[}i{]}

โดยมีค่าใช้จ่าย c{[}i{]} โดยที่คุณสมบัติของพนักงาน a{[}i{]} จะต้องมากกว่าพนักงาน

b{[}i{]} เสมอ (q{[}a{[}i{]}{]} \textgreater{} q{[}b{[}i{]}{]})



\subsection*{Output}
บรรทัดเดียว --- ค่าใช้จ่ายที่ต่ำที่สุดในการสร้างโครงสร้างลำดับชั้น ``หัวหน้างาน-ลูกน้อง'' นี้

หรือ -1 หากไม่สามารถสร้างโครงสร้างนี้ได้\\



\end{document}
